\chapter*{ABSTRACT}
\addcontentsline{toc}{chapter}{Abstract}

Long Term Evolution networks execute handovers through Event A3, a rule that fires only after a neighbour cell has already become stronger than the serving cell by a configured offset and has remained stronger for a full time-to-trigger. The rule is reactive by construction: it issues a response, not a warning. This thesis asks whether the next handover can be forecast from the measurements a handset already reports, how far ahead, with what guarantee --- and, as it turned out, what has to be verified before any such number can be believed.

Four drive-test campaigns were conducted on a live commercial LTE network in Dhaka and Gazipur, yielding 4 continuous logging sessions, 10,260 samples at one hertz and 957 handover commands decoded from RRC signalling (yielding 938 quality-controlled commands inside retained session blocks, and 934--937 horizon-evaluable events under lead-window boundaries). An alignment audit of the one-hertz export is reported first, because it governs every result that follows: in 76\% of handovers the export row stamped before the handover command already carries the post-handover serving cell, so a row summarises the second that follows its own timestamp. Every feature is therefore lagged by one row, which moves the one-second AUPRC of the same model from 0.600 to 0.179; an arm-wise decomposition attributes 96\% of that fall to the serving-cell radio scalars (where leak elimination and one-second feature ageing are confounded) and negligible impact to lagging kinematics. The defect is consistent with an end-of-second aggregation mechanism: a row written at the end of its second carries the target cell whenever the handover completes before that second closes, an empirical hypothesis consistent with an in-sample diagnostic check on the 15 handovers here whose execution crosses the boundary (0/15; 95\% Wilson score interval [0.0\%, 20.4\%], while the 211 residual cases are consistent with an additional 1-row exporter refresh delay). It replicates across instruments on 4,712 handovers from an external multi-carrier XCAL corpus spanning 3 operators [40] (93\% contaminated under an end-of-second rule, 4\% under a start-of-second rule on the same events) and is absent from a public log whose rows are instantaneous samples [37]. Read empirically, the transition cliff characterises the sub-second logging and periodic export commit latency around 50--75 ms from a one-hertz export. The task is posed as discrete-time survival: one model estimates per-interval hazards and horizon probabilities are products of hazards, so ordering holds by construction for any learner.

Under leave-one-campaign-out evaluation --- the protocol whose units represent distinct physical sessions here --- gradient-boosted trees reach pooled AUPRC 0.179 at one second across the four campaigns (empirical campaign spread [0.132, 0.233]; pooled 180~s block-bootstrap 95\% CI [0.159, 0.204]) against a 6.8\% prevalence floor (2.6x lift, AUROC 0.777, median lead time 0.47~s) and AUPRC 0.475 at five seconds (1.8x, campaign spread [0.412, 0.548]), against 0.116 for the implemented A3 signal-margin proxy scored as a predictor. Splitting rows at random instead inflates the same number by 48 \%, and splitting by 180-second chunks of one session --- the protocol this work used before the audit --- still inflates it by 13 \%. The hazard formulation yields zero ordering violations against 40.2\% of rows for independent per-horizon classifiers, a defect that a cumulative-maximum projection also removes at no cost, so the formulation is defended on needing neither a held-out calibration split nor five separate fits rather than on coherence alone. Scored as a warning rather than as a ranking, the system detects 34\% of handover commands at the one-second horizon at 179 false alarms per hour with a median lead of 0.47 s, and 42\% at five seconds at 70 per hour with a median lead of 3.31 s; no operating point measured here is deployable for target-cell preparation, and the contribution is the measurement framework rather than a deployable predictor. Conformal risk control is investigated as a calibration diagnostic rather than an operating guarantee: distribution-free finite-sample guarantees cannot be certified across campaigns because exchangeability is violated by macro-environmental covariate shift (with 12 rotations exhibiting pseudoreplication over $n=4$ physical sessions), and 35.8\% of handover commands lack an eligible prediction row at 1~s under the enforced 2~s post-handover blanking floor. As an empirical calibration diagnostic, the pooled mean miss loss across rotations was 0.153 at nominal $\alpha=0.20$.

Two properties of the deployed network are corrected against the earlier draft of this work. Weighted by handovers rather than by configurations, 75\% of A3-triggered handovers fire under a positive offset, so the operator's negative offsets belong to measurement configurations that rarely lead to mobility. The A3 conversion rate is shown to have no unit-free value: 63\% of reports, 30\% of trigger episodes under a loose rule that lets overlapping episodes claim the same command, and 73\% once each command is matched to one episode. The separation between mobility and monitoring configurations survives all three and is widest under the strictest: 37\% declined for the dominant mobility profile against 95\% for the dominant monitoring profile. The contribution is a signalling-grounded, leakage-audited framework that converts measured handover events into coherent multi-horizon probabilities with a stated operating cost, together with a reproducible demonstration of how much of an apparently strong result on drive-test data can be timestamp alignment.

\vspace{1em}\noindent\textbf{Keywords:} LTE, handover prediction, Event A3, RRC signalling, discrete-time survival analysis, data leakage, evaluation protocol, conformal risk control, calibration, drive test.

